% Options for packages loaded elsewhere
\PassOptionsToPackage{unicode}{hyperref}
\PassOptionsToPackage{hyphens}{url}
%
\documentclass[
]{article}
\usepackage{amsmath,amssymb}
\usepackage{iftex}
\ifPDFTeX
  \usepackage[T1]{fontenc}
  \usepackage[utf8]{inputenc}
  \usepackage{textcomp} % provide euro and other symbols
\else % if luatex or xetex
  \usepackage{unicode-math} % this also loads fontspec
  \defaultfontfeatures{Scale=MatchLowercase}
  \defaultfontfeatures[\rmfamily]{Ligatures=TeX,Scale=1}
\fi
\usepackage{lmodern}
\ifPDFTeX\else
  % xetex/luatex font selection
\fi
% Explicit Unicode mapping for the pandoc-generated body under pdflatex
% (utf8 inputenc does not define Greek / math symbols by default).
\ifPDFTeX
\usepackage{newunicodechar}
\newunicodechar{Λ}{\ensuremath{\Lambda}}
\newunicodechar{Σ}{\ensuremath{\Sigma}}
\newunicodechar{Φ}{\ensuremath{\Phi}}
\newunicodechar{β}{\ensuremath{\beta}}
\newunicodechar{λ}{\ensuremath{\lambda}}
\newunicodechar{μ}{\ensuremath{\mu}}
\newunicodechar{ρ}{\ensuremath{\rho}}
\newunicodechar{τ}{\ensuremath{\tau}}
\newunicodechar{—}{---}
\newunicodechar{−}{\textminus}
\newunicodechar{≅}{\ensuremath{\cong}}
\newunicodechar{≈}{\ensuremath{\approx}}
\newunicodechar{⊗}{\ensuremath{\otimes}}
\fi
% Use upquote if available, for straight quotes in verbatim environments
\IfFileExists{upquote.sty}{\usepackage{upquote}}{}
\IfFileExists{microtype.sty}{% use microtype if available
  \usepackage[]{microtype}
  \UseMicrotypeSet[protrusion]{basicmath} % disable protrusion for tt fonts
}{}
\makeatletter
\@ifundefined{KOMAClassName}{% if non-KOMA class
  \IfFileExists{parskip.sty}{%
    \usepackage{parskip}
  }{% else
    \setlength{\parindent}{0pt}
    \setlength{\parskip}{6pt plus 2pt minus 1pt}}
}{% if KOMA class
  \KOMAoptions{parskip=half}}
\makeatother
\usepackage{xcolor}
\setlength{\emergencystretch}{3em} % prevent overfull lines
\providecommand{\tightlist}{%
  \setlength{\itemsep}{0pt}\setlength{\parskip}{0pt}}
\setcounter{secnumdepth}{-\maxdimen} % remove section numbering
\ifLuaTeX
  \usepackage{selnolig}  % disable illegal ligatures
\fi
\usepackage{bookmark}
\IfFileExists{xurl.sty}{\usepackage{xurl}}{} % add URL line breaks if available
\urlstyle{same}
\hypersetup{
  hidelinks,
  pdfcreator={LaTeX via pandoc}}

\author{}
\date{}

\begin{document}

\section{Relational Emergence Model
(REM)}\label{relational-emergence-model-rem}

\subsection{Core Specification v2.0 (Frozen 2026-08-11)}\label{core-specification-v2.0}

This version supersedes v1.1. Changes: open-system general form
$\Phi(F;\lambda,\rho,\mathcal L,\tau) = I_\rho(F) - \lambda C_{\mathrm{dyn}}(F;\rho,\mathcal L,\tau)$
with $[\lambda] = T$;
$C_H^{\mathrm{closed}}$ downgraded to a closed-system structural surrogate;
Monotonic Tradeoff Theorem extended to signed dynamical functionals.

\section{0. Core Statement}\label{core-statement}

REM treats tensor-product structure of Hilbert space as a variable
structure selected via competing informational and dynamical
constraints.

REM is a structural selection framework, not a modification of quantum
dynamics.

\section{1. Three-Layer Architecture}\label{three-layer-architecture}

\subsection{Layer 0 --- Pre-relational
description}\label{layer-0-pre-relational-description}

Hilbert space H is considered without specifying a preferred
tensor-product decomposition.

Layer 0 does not represent a physical substrate. It denotes absence of
articulated relational structure.

\subsection{Layer 1 --- Relational
articulation}\label{layer-1-relational-articulation}

Differentiation introduces an explicit tensor structure:

\[
\mathcal D : \mathcal H \to \mathcal H_S \otimes \mathcal H_O
\]

A quantum state admits Schmidt representation:

\[
\mathcal D(|\Psi\rangle) \to \sum_i c_i |s_i\rangle \otimes |o_i\rangle
\]

Relational structure becomes explicit.

\subsection{Layer 2 --- Phenomenal
stabilization}\label{layer-2-phenomenal-stabilization}

Environmental decoherence stabilizes articulated correlations:

\[
\mathcal E(\rho) = \sum_i p_i |s_i\rangle\langle s_i| \otimes |e_i\rangle\langle e_i|
\]

Ordering relation:

\[
\mathcal D \to \mathcal E
\]

\section{2. Differentiation map}\label{differentiation-map}

D is a structural selection map.

D is not a dynamical law and is not a unitary operator.

D is defined implicitly through a selection problem over candidate
factorizations.

Variational principles provide tractable models for D but do not exhaust
its possible formalizations.

\section{3. Factorization space}\label{factorization-space}

Candidate tensor decompositions:

F = \{F\_k\}

Finite-dimensional approximation:

F ≅ U(N)/(U(n\_A) ⊗ U(n\_B))

This representation is heuristic and applies primarily in
finite-dimensional Hilbert spaces.

\section{4. Variational ansatz}\label{variational-ansatz}

General form (open-system REM, 2026-08-11):

Φ(F; λ, ρ, L, τ) = I\_ρ(F) − λ C\_dyn(F; ρ, L, τ)

I\_ρ is the mutual information of the (possibly mixed) state ρ across F.

Closed-system approximation (REM\_lambda v4 and earlier):

\[
\Phi_{\mathrm{closed}}(F) = I(F) - \lambda C_H^{\mathrm{closed}}(F)
\]

Higher-order generalization:

Φ(F) = I − λ C\_dyn + Σ μ\_i Φ\_i

Optimal factorization:

F* = argmax\_F Φ(F)

The functional provides an effective model for relational selection.

\section{5. Informational functional}\label{informational-functional}

Φ\_S = I(A:B)

I(A:B) = S(ρ\_A) + S(ρ\_B) − S(ρ\_AB)

Von Neumann entropy:

S(ρ) = −Tr(ρ log ρ)

Units may be nats or bits depending on logarithm base.

\section{6. Dynamical functional}\label{dynamical-functional}

C\_dyn measures dynamical (open-system) coupling across candidate
partitions. It is environment-dependent: C\_dyn = C\_dyn(F; ρ, L, τ).
Formal name: signed dynamical functional — the SIGN of C\_dyn matters
(negative values reward rather than penalize a structure; "cost" is used
only informally).

First operational realization (2026-08-11, Phase C):

\[
C_\Gamma^{(0)}(F;\rho,\mathcal L)
= \Gamma_F^{\mathrm{exact}}(0)
= -\frac{\mathrm{Re}\langle Q_F\rho,\,Q_F\mathcal L(\rho)\rangle_{\mathrm{HS}}}{|Q_F\rho|_2^2}
\]

with $Q_F = I - D_F$ ($D_F$: dephasing projection onto the initial Schmidt
basis of $F$), $\mathcal L$ the Lindblad Liouvillian of the fixed
environment, and $\langle A,B\rangle_{\mathrm{HS}} = \mathrm{Tr}(A^\dagger B)$.
This is the initial structural decoherence rate.
For pure states it is independent of the Schmidt null-subspace completion.
A finite-time variant $C_{\mathrm{dyn}}^{(\tau)} = -(1/\tau)\log(C_F(\tau)/C_F(0))$
is the natural extension when the timescale $\tau$ matters.

SIGN: $C_\Gamma^{(0)}$ is a SIGNED rate. Random-state audit (2026-08-11, 400
configs) found negative values (1.25\%, min $\approx -1.6$) — initial coherence
increase is possible; the quantity is a signed structural decoherence
rate, not a non-negative cost, and must not be forced through $\max(0,\cdot)$.

Closed-system surrogate (REM\_lambda v4 convention):

\[
C_H^{\mathrm{closed}}(F) = \mathrm{Tr}(\rho H_{\partial F}^2)
\]

$H_{\partial F}$ represents interaction terms spanning the factorization
boundary. $C_H^{\mathrm{closed}}$ is the boundary-energy second moment, not the
variance: $C_H^{\mathrm{closed}} = \langle H_{\partial F}^2\rangle \neq
\mathrm{Var}(H_{\partial F}) = \langle H_{\partial F}^2\rangle -
\langle H_{\partial F}\rangle^2$.
It is a closed-system structural surrogate, not an open-system stability
cost (C0--C2, 2026-08-11).

Dimensions (canonical, frozen 2026-08-11): $[C_{\mathrm{dyn}}] = T^{-1}$
and $[\lambda] = T$, so $\lambda C_{\mathrm{dyn}}$ is dimensionless.
$\lambda$ is the characteristic timescale comparing informational
articulation and dynamical stability. The finite-time variant
$C_{\mathrm{dyn}}^{(\tau)}$ also has $[C_{\mathrm{dyn}}^{(\tau)}] = T^{-1}$,
so the same dimensionful $\lambda$ applies. The normalized representation
$\tilde C_{\mathrm{dyn}} = \tau_0 C_{\mathrm{dyn}}$,
$\tilde\lambda = \lambda/\tau_0$ (with
$\Phi = I - \tilde\lambda \tilde C_{\mathrm{dyn}}$) is mathematically
equivalent and kept only as an optional representation; the canonical
form introduces no new reference time $\tau_0$.

The earlier linear ansatz Φ\_H = −Tr(ρ H\_boundary) is superseded.

Interpretation:

dynamical self-containment measure.

\section{7. Interpretation of λ}\label{interpretation-of-ux3bb}

Primary meaning:

relative weighting between informational articulation and dynamical
self-containment.

Secondary interpretations:

effective thermodynamic analogy:

Φ ≈ S − βE

emergent scale parameter:

λ = λ(Λ)

timescale relation:

λ \textasciitilde{} τ\_D / τ\_dyn

λ is an effective parameter, not a universal constant.

\section{8. Crossover structure}\label{crossover-structure}

λ* defined by:

Φ(F1, λ\emph{) = Φ(F2, λ})

λ* is model dependent.

REM predicts existence of crossover regimes, not universal λ* values.

\section{8a. Monotonic Tradeoff Theorem}\label{monotonic-tradeoff-theorem}

Let $\Phi_\lambda(F) = S(F) - \lambda C(F)$ with $\lambda \geq 0$, and let
$F^*(\lambda)$ denote a global maximizer of $\Phi_\lambda$.

Theorem (Monotonic Tradeoff). If $\lambda_2 > \lambda_1 \geq 0$, then

\[
C(F^*(\lambda_2)) \leq C(F^*(\lambda_1)), \qquad
S(F^*(\lambda_2)) \leq S(F^*(\lambda_1)).
\]

Proof. Let $F_1 = F^*(\lambda_1)$ and $F_2 = F^*(\lambda_2)$. Global
optimality gives

\[
S_1 - \lambda_1 C_1 \geq S_2 - \lambda_1 C_2, \qquad
S_2 - \lambda_2 C_2 \geq S_1 - \lambda_2 C_1.
\]

Adding the two inequalities yields

\[
(\lambda_2 - \lambda_1)(C_1 - C_2) \geq 0
\]

so $C_2 \leq C_1$. Substituting into the first inequality,

\[
S_1 - S_2 \geq \lambda_1(C_1 - C_2) \geq 0
\]

so $S_2 \leq S_1$. $\square$

The theorem holds for any non-negative cost functional $C$ of the
factorization --- in particular both the second moment
$C_H = \langle H_{\partial}^2\rangle$ and the variance
$\mathrm{Var}(H_{\partial})$ --- and requires only that $F^*(\lambda)$
is a global maximizer.

In the open-system generalization the theorem applies unchanged for a
FIXED environment $\mathcal L$, state $\rho$, and timescale $\tau$:
$C_{\mathrm{dyn}}(F;\rho,\mathcal L,\tau)$ is a well-defined functional
of $F$ for each fixed $(\rho,\mathcal L,\tau)$, so the monotonicity
$C(F^*(\lambda_2)) \leq C(F^*(\lambda_1))$ and
$S(F^*(\lambda_2)) \leq S(F^*(\lambda_1))$ holds per environment (the
proof requires only $\lambda \geq 0$; the cost $C$ may be signed, as in
the $C_\Gamma^{(0)}$ realization).

\section{9. Role of toy models}\label{role-of-toy-models}

Small systems illustrate nontrivial structure in factorization
landscape.

Toy models serve as existence proofs.

They do not claim empirical completeness.

\section{10. Relation to prior work}\label{relation-to-prior-work}

REM is situated within ongoing discussions concerning the non-uniqueness of subsystem structure in quantum theory.

Zanardi, Lidar, and Whaley (2004) demonstrated that tensor product structure is not uniquely determined by the Hilbert space alone, providing an algebraic classification of admissible subsystem decompositions. However, no physical selection principle was proposed for choosing among admissible factorizations.

Cotler, Penington, and Ranard (2019) showed that preferred factorizations may arise through optimization procedures, but the associated objective function was not interpreted as expressing a competition between informational and dynamical constraints.

Carroll and Singh (2021) proposed a quantum mereological framework in which subsystem structure is constrained by dynamical considerations. Their construction emphasizes quasi-classical persistence but does not formulate the selection problem as a variational competition between informational articulation and dynamical self-containment.

Donnelly and Freidel (2016) demonstrated that subsystem structure in gauge theories is nontrivial due to boundary degrees of freedom, indicating that factorization is not generically available as a primitive notion.

Rovelli (1996) introduced relational quantum mechanics, in which physical states are defined relative to interactions between systems. However, the theory does not specify how the relational partition itself is determined.

Zurek (2003) showed that environment-induced decoherence selects stable pointer observables, thereby explaining classical phenomenology once subsystem structure is assumed.

REM contributes a complementary structural proposal:

the tensor product structure itself is treated as a variable object selected through a variational competition between informational articulation and dynamical self-containment.

Within this framework, crossover behaviour between competing structural criteria is characterized by an effective parameter λ, producing regime-dependent preferred factorizations.

REM therefore addresses a structural layer prior to relational state assignment and decoherence-based stabilization.

\section{11. Scope of theory}\label{scope-of-theory}

REM addresses structural selection problem.

REM does not introduce new dynamical laws, hidden variables, or
modifications to quantum mechanics.

REM provides structural constraints on subsystem articulation.

\section{12. Terminology}\label{terminology}

differentiation: selection of relational tensor structure

factorization: tensor product decomposition

articulation: explicit correlation structure

self-containment: weak dynamical coupling across partition

crossover: change in preferred factorization regime

\section{13. Extension directions}\label{extension-directions}

multi-qubit scaling

tensor-network articulation

continuous Hilbert spaces

renormalization structure

quantum channel formulations of differentiation

\section{14. Status of theory}\label{status-of-theory}

REM is a structured foundational proposal.

It provides a mathematically formulated research program for relational
structure selection.

Future work will determine empirical and mathematical generality.

\end{document}
